\documentclass[a4paper,12pt]{article}

\input{../../Preambulo.tex}
\newcommand{\assunto}{FUNÇÕES AFINS}
\newcommand{\atividade}{lista} % prova ou lista
\newcommand{\turma}{ELETROELETRÔNICA 1.\textordmasculine\ ANO}
\newcommand{\numProva}{A}
\newcommand{\professor}{Igor Martins Silva}
\newcommand{\bimestre}{3}
\newcommand{\ano}{2026}
\newcommand{\mesTexto}{agosto}
\newcommand{\mesNum}{08}
\newcommand{\dia}{27}
\newcommand{\dataTexto}{\dia\ de \mesTexto\ de \ano}
\newcommand{\dataNun}{\dia/\mesNum/\ano}

\newcommand{\titulo}{}

\ifdefstring{\atividade}{prova}
  {\renewcommand{\titulo}{Avaliação de Matemática}}
{
\ifdefstring{\atividade}{lista}
  {\renewcommand{\titulo}{Lista de exercícios de Matemática}}
{
\ifdefstring{\atividade}{as}
  {\renewcommand{\titulo}{Avaliação Somativa}}
{
\ifdefstring{\atividade}{ad}
  {\renewcommand{\titulo}{Avaliação Diagnóstica}}
  {\renewcommand{\titulo}{Atividade de Matemática}}
}}}

\newcommand{\emtipo}{%
  \ifdefstring{\tipo}{individual}{\tipo}{em \tipo}%
}

\edef\pathCapa{\currfiledir}

\newcommand{\subtitulobase}[3]{%
    % #1 = título da 1ª coluna
    % #2 = conteúdo da 1ª coluna
    % #3 = mostra número da prova? (vazio ou algo)
    
    \noindent
    \begin{minipage}{2.8cm}
        \includegraphics[width=2.8cm]{\pathCapa Logo_Cefet.png}
    \end{minipage}
    \hfill
    \begin{minipage}{\dimexpr\textwidth-6.2cm\relax}
        \begin{center}
        CEFET-MG -- Campus Contagem \\[3pt]
        \titulo\ -- \bimestre.\textordmasculine\ bimestre de \ano \\[3pt]
        \professor
    \end{center}
    \end{minipage}
    \hfill
    \begin{minipage}{3.2cm}
        \includegraphics[width=3.2cm]{\pathCapa Brasao_Brasil.png}
    \end{minipage}
    
    \begin{center}
        \vspace{15pt}
        \begin{tabular}{|C{210pt}|C{70pt}|C{210pt}|}
            \hline
            \textbf{#1} &
            \textbf{DATA} &
            \textbf{TURMA} \\
            \hline
            #2 & \dataNun & \turma \\
            \hline
        \end{tabular}
    \end{center}
    
    #3
    
    \vspace{15pt}
}

\newcommand{\subtitulolis}{%
  \subtitulobase{ASSUNTO}{\assunto}{}%
}

\newcommand{\subtitulopr}{%
  \subtitulobase{ASSUNTO}{\assunto}{\boxed{\numProva}}%
}

\newcommand{\subtituloas}{%
  \subtitulobase{DISCIPLINA}{MATEMÁTICA}{\boxed{\numProva}}%
}

\newcommand{\subtituload}{%
  \subtitulobase{DISCIPLINA}{MATEMÁTICA}{}%
}

\newcommand{\valorquestao}{}

\newcommand{\definevalorquestao}{%
    \ifdefstring{\atividade}{lista}{%
        % Não faz nada para lista
    }{%
        \ifdefstring{\modo}{objetivo}{%
            \renewcommand{\valorquestao}{\ValorQObj ponto}
        }{%
            \renewcommand{\valorquestao}{\ValorQDisc pontos}
        }%
    }%
}

\newenvironment{exercicioBanco}[1][]{%
    \ifdefstring{\atividade}{lista}{%
        \begin{exercicio}%
    }{%
        \begin{exercicio}[#1]%
    }%
}{%
    \end{exercicio}
}

\ifdefstring{\atividade}{lista}{
    \pagecolor{background-color}
}{
    
}



\hypersetup{
    pdfauthor={\professor},
    pdftitle={\titulo \ -- \assunto \ -- \turma},
}

\title{\titulo \ -- \assunto \ -- \turma}
\author{\professor}
\date{\data}

\graphicspath{
    {../../FuncoesAfins/}
}

%************* INÍCIO DO DOCUMENTO *************
\begin{document}
    \subtitulolis
        
    \phantomsection
    \addcontentsline{toc}{section}{Exercício 1}
    %%%%%%%%%%%%%%%%%%%%%%%%%%%%%%%%%
%
%       EXERCÍCIO
%
%%%%%%%%%%%%%%%%%%%%%%%%%%%%%%%%%

\ifdefstring{\atividade}{prova}{%
    \ifdefstring{\modo}{objetivo}{%
        \renewcommand{\valorquestao}{\ValorQObj\ ponto}
    }{%
        \renewcommand{\valorquestao}{\ValorQDisc\ pontos}
    }%
}%

\begin{exercicioBanco}[\valorquestao]
Determine quais das funções a seguir são funções afins.
%
\begin{center}
    \begin{tabularx}{\textwidth}{XXXX}
        (I) \(f(x) = 0\). & (II) \(f(x) = x^{2}\). & (III) \(f(x) = \pi x - \sqrt{2}\). & (IV) \(f(x) = 2^{x}\).  \\[10pt]
        (V) \(f(x) = x\). & (VI) \(f(x) = -\dfrac{1}{2} - 2x\). & (VII) \(f(x) = \dfrac{1}{x}\). & (VIII) \(f(x) = \sqrt{x}\).
    \end{tabularx}
\end{center}

% Define as alternativas
\newcommand{\alternativas}{%
    \begin{center}
        \begin{tabularx}{\textwidth}{XXXXX}
            (a) I, III, V, VI. &
            (b) V, VI. &
            (c) V, VI, VII. &
            (d) Todas. &
            (e) Todas, exceto IV.
        \end{tabularx}
    \end{center}
}

% Define a resposta correta
\newcommand{\resposta}{A}

% Lógica condicional para exibição
\ifdefstring{\atividade}{lista}{%
    \alternativas
    \vspace{0.5em}
    
    \noindent\textbf{Resposta:} letra \textbf{\resposta}.
}{%
    \ifdefstring{\modo}{objetivo}{%
        \alternativas
    }{%
        % modo = discursiva → não mostra alternativas
    }
}
\end{exercicioBanco}


    \noindent\rule{\linewidth}{1pt}  % linha horizontal fina
    
    \phantomsection
    \addcontentsline{toc}{section}{Exercício 2}
    \input{../FA_Exercicio2_1}
    \noindent\rule{\linewidth}{1pt}  % linha horizontal fina
    
    \phantomsection
    \addcontentsline{toc}{section}{Exercício 3}
    %%%%%%%%%%%%%%%%%%%%%%%%%%%%%%%%%
%
%       EXERCÍCIO
%
%%%%%%%%%%%%%%%%%%%%%%%%%%%%%%%%%

\ifdefstring{\atividade}{prova}{%
    \ifdefstring{\modo}{objetivo}{%
        \renewcommand{\valorquestao}{\ValorQObj\ ponto}
    }{%
        \renewcommand{\valorquestao}{\ValorQDisc\ pontos}
    }%
}%

\begin{exercicioBanco}[\valorquestao]
Seja \(f\) uma função afim tal que \(f(0) = -1\) e \(f(2) = 0\). Determine a lei que representa essa função.

% Define as alternativas
\newcommand{\alternativas}{%
    \begin{center}
        \begin{tabularx}{\textwidth}{XXX}
            (a) \(f(x) = \dfrac{1}{2}x + 1\). &
            (b) \(f(x) = -\dfrac{1}{2}x - 1\). &
            (c) \(f(x) = \dfrac{1}{2}x - 1\). \\[10pt]
            (d) \(f(x) = x - 1\). &
            (e) \(f(x) = \dfrac{1}{2}x - 2\).
        \end{tabularx}
    \end{center}
}

% Resposta correta
\newcommand{\resposta}{C}

% Lógica condicional para exibição
\ifdefstring{\atividade}{lista}{%
    \alternativas
    \vspace{0.5em}
    
    \noindent\textbf{Resposta:} letra \textbf{\resposta}.
}{%
    \ifdefstring{\modo}{objetivo}{%
        \alternativas
    }{%
        % modo = discursiva → não mostra alternativas
    }
}
\end{exercicioBanco}


    \noindent\rule{\linewidth}{1pt}  % linha horizontal fina
    
    \phantomsection
    \addcontentsline{toc}{section}{Exercício 4}
    %%%%%%%%%%%%%%%%%%%%%%%%%%%%%%%%%
%
%       EXERCÍCIO
%
%%%%%%%%%%%%%%%%%%%%%%%%%%%%%%%%%

\ifdefstring{\atividade}{prova}{%
    \ifdefstring{\modo}{objetivo}{%
        \renewcommand{\valorquestao}{\ValorQObj\ ponto}
    }{%
        \renewcommand{\valorquestao}{\ValorQDisc\ pontos}
    }%
}%

\begin{exercicioBanco}[\valorquestao]
Seja a função afim 
%
\[
    \begin{array}{rcl}
        f:\ \bbR & \to & \bbR \\
        x & \mapsto & 2x.
    \end{array}
\]
%
Calcule o valor de \(f(2) + f(3) - f(1)\).

% Define as alternativas
\newcommand{\alternativas}{%
    \begin{center}
        \begin{tabularx}{\textwidth}{XXXXX}
            (a) \(12\). &
            (b) \(10\). &
            (c) \(4\). &
            (d) \(8\). &
            (e) \(6\).
        \end{tabularx}
    \end{center}
}

% Define a resposta correta
\newcommand{\resposta}{D}

% Lógica condicional para exibição
\ifdefstring{\atividade}{lista}{%
    \alternativas
    \vspace{0.5em}
    
    \noindent\textbf{Resposta:} letra \textbf{\resposta}.
}{%
    \ifdefstring{\modo}{objetivo}{%
        \alternativas
    }{%
        % modo = discursiva → não mostra alternativas
    }
}
\end{exercicioBanco}


    \noindent\rule{\linewidth}{1pt}  % linha horizontal fina
    
    \phantomsection
    \addcontentsline{toc}{section}{Exercício 5}
    %%%%%%%%%%%%%%%%%%%%%%%%%%%%%%%%%
%
%       EXERCÍCIO
%
%%%%%%%%%%%%%%%%%%%%%%%%%%%%%%%%%

\ifdefstring{\atividade}{prova}{%
    \ifdefstring{\modo}{objetivo}{%
        \renewcommand{\valorquestao}{\ValorQObj\ ponto}
    }{%
        \renewcommand{\valorquestao}{\ValorQDisc\ pontos}
    }%
}%

\begin{exercicioBanco}[\valorquestao]
Seja a função afim \(f: \bbR \to \bbR\), definida por \(f(x) = -x + \frac{1}{2}\). Determine seus coeficientes angular e linear, indicando-os nessa ordem.

% Define as alternativas
\newcommand{\alternativas}{%
    \begin{center}
        \begin{tabularx}{\textwidth}{XXXXX}
            (a) \(1\) e \(-\dfrac{1}{2}\). &
            (b) \(1\) e \(\dfrac{1}{2}\). &
            (c) \(-1\) e \(-\dfrac{1}{2}\). &
            (d) \(\dfrac{1}{2}\) e \(-1\). &
            (e) \(-1\) e \(\dfrac{1}{2}\).
        \end{tabularx}
    \end{center}
}

% Define a resposta correta
\newcommand{\resposta}{E}

% Lógica condicional para exibição
\ifdefstring{\atividade}{lista}{%
    \alternativas
    \vspace{0.5em}
    
    \noindent\textbf{Resposta:} letra \textbf{\resposta}.
}{%
    \ifdefstring{\modo}{objetivo}{%
        \alternativas
    }{%
        % modo = discursiva → não mostra alternativas
    }
}
\end{exercicioBanco}


    \noindent\rule{\linewidth}{1pt}  % linha horizontal fina
    
    \phantomsection
    \addcontentsline{toc}{section}{Exercício 6}
    %%%%%%%%%%%%%%%%%%%%%%%%%%%%%%%%%
%
%       EXERCÍCIO
%
%%%%%%%%%%%%%%%%%%%%%%%%%%%%%%%%%

\ifdefstring{\atividade}{prova}{%
    \ifdefstring{\modo}{objetivo}{%
        \renewcommand{\valorquestao}{\ValorQObj\ ponto}
    }{%
        \renewcommand{\valorquestao}{\ValorQDisc\ pontos}
    }%
}%

\begin{exercicioBanco}[\valorquestao]
Considere a função afim cujo gráfico está desenhado abaixo.
%

%
\begin{center}
    \begin{tikzpicture}[scale=0.6]
        % Eixos
        \draw[->] (-3,0) -- (2.5,0) node[right] {\(x\)};
        \draw[->] (0,-2) -- (0,2) node[above] {\(y\)};
              
        % Gráfico
        \draw[domain=-2.5:1.5, smooth, variable=\x, blue, thick] 
        plot ({\x}, {\x+1});
              
        % Pontos escolhidos
        \draw[dashed] (-2,0) -- (-2,-1) -- (0,-1);
        \draw[dashed] (1,0) -- (1,2) -- (0,2);
             
        % Marcando os pontos
        \filldraw (-2,-1) circle (1pt) node[left] {\small\((-2,-1)\)};
        \filldraw (1,2) circle (1pt) node[right] {\small\((1,2)\)};
    \end{tikzpicture}
\end{center}
%
Determine o zero dessa função.

% Define as alternativas
\newcommand{\alternativas}{%
    \begin{center}
        \begin{tabularx}{\textwidth}{XXXXX}
            (a) \(-2\). &
            (b) \(-1\). &
            (c) \(0\). &
            (d) \(1\). &
            (e) \(2\).
        \end{tabularx}
    \end{center}
}

% Define a resposta correta
\newcommand{\resposta}{B}

% Lógica condicional para exibição
\ifdefstring{\atividade}{lista}{%
    \alternativas
    \vspace{0.5em}
    
    \noindent\textbf{Resposta:} letra \textbf{\resposta}.
}{%
    \ifdefstring{\modo}{objetivo}{%
        \alternativas
    }{%
        % modo = discursiva → não mostra alternativas
    }
}
\end{exercicioBanco}


    \noindent\rule{\linewidth}{1pt}  % linha horizontal fina
    
    \phantomsection
    \addcontentsline{toc}{section}{Exercício 7}
    %%%%%%%%%%%%%%%%%%%%%%%%%%%%%%%%%
%
%       EXERCÍCIO
%
%%%%%%%%%%%%%%%%%%%%%%%%%%%%%%%%%

\ifdefstring{\atividade}{prova}{%
    \ifdefstring{\modo}{objetivo}{%
        \renewcommand{\valorquestao}{\ValorQObj\ ponto}
    }{%
        \renewcommand{\valorquestao}{\ValorQDisc\ pontos}
    }%
}%

\begin{exercicioBanco}[\valorquestao]
Determine o gráfico da função afim \(f: \bbR \to \bbR\), dada por \(f(x) = 3x + 2\).

% Define as alternativas
\newcommand{\alternativas}{%
\begin{center}
    \begin{tabular}{C{20pt}C{140pt}C{20pt}C{140pt}C{20pt}C{140pt}}
        (a) 
        &
        \begin{tikzpicture}[scale=0.5]
            % Eixos
            \draw[->] (-3.0,0) -- (3.0,0) node[right] {\(x\)};
            \draw[->] (0,-3.0) -- (0,3.0) node[above] {\(y\)};

            % Gráfico da função f(x) = |x+1| - 2
            \draw[domain=-2.8:2.8, thick, blue] plot (\x, {\x});
        \end{tikzpicture}
        &
        (b) 
        &
        \begin{tikzpicture}[scale=0.5]
            % Eixos
            \draw[->] (-3.0,0) -- (3.0,0) node[right] {\(x\)};
            \draw[->] (0,-3.0) -- (0,3.0) node[above] {\(y\)};

            % Gráfico da função
            \draw[domain=-1.5:0.5, thick, blue] plot (\x, {3*(\x) + 2});

            % Marcações da raiz
            \node[above left] at (-0.67,0) {\small\(-\frac{2}{3}\)};
            \node at (-0.67,0) {\small\(\bullet\)};

            % Marcação do ponto b
            \node[right] at (0,2) {\small\(2\)};
            \node at (0,2) {\small\(\bullet\)};
        \end{tikzpicture}
        &
        (c)
        &
        \begin{tikzpicture}[scale=0.5]
            % Eixos
            \draw[->] (-3.0,0) -- (3.0,0) node[right] {\(x\)};
            \draw[->] (0,-3.0) -- (0,3.0) node[above] {\(y\)};

            % Gráfico da função
            \draw[domain=-0.5:1.5, thick, blue] plot (\x, {3*(\x) - 2});

            % Marcações da raiz
            \node[below right] at (0.67,0) {\small\(\frac{2}{3}\)};
            \node at (0.67,0) {\small\(\bullet\)};

            % Marcação do ponto b
            \node[left] at (0,-2) {\small\(-2\)};
            \node at (0,-2) {\small\(\bullet\)};
        \end{tikzpicture}
        \\
        (d)
        &
        \begin{tikzpicture}[scale=0.5]
            % Eixos
            \draw[->] (-3.0,0) -- (3.0,0) node[right] {\(x\)};
            \draw[->] (0,-3.0) -- (0,3.0) node[above] {\(y\)};

            % Gráfico da função
            \draw[domain=-2.7:1.5, thick, blue] plot (\x, {\x + 2});

            % Marcações da raiz
            \node[above left] at (-2,0) {\small\(-2\)};
            \node at (-2,0) {\small\(\bullet\)};

            % Marcação do ponto b
            \node[right] at (0,2) {\small\(2\)};
            \node at (0,2) {\small\(\bullet\)};
        \end{tikzpicture}
        &
        (e) 
        &
        \begin{tikzpicture}[scale=0.5]
            % Eixos
            \draw[->] (-3.0,0) -- (3.0,0) node[right] {\(x\)};
            \draw[->] (0,-3.0) -- (0,3.0) node[above] {\(y\)};

            % Gráfico da função
            \draw[domain=-0.5:1.5, thick, blue] plot (\x, {-3*(\x) + 2});

            % Marcações da raiz
            \node[above right] at (0.67,0) {\small\(\frac{2}{3}\)};
            \node at (0.67,0) {\small\(\bullet\)};

            % Marcação do ponto b
            \node[left] at (0,2) {\small\(2\)};
            \node at (0,2) {\small\(\bullet\)};
        \end{tikzpicture}
        &&
    \end{tabular}
\end{center}
}

% Define a resposta correta
\newcommand{\resposta}{B}

% Lógica condicional para exibição
\ifdefstring{\atividade}{lista}{%
    \alternativas
    \vspace{0.5em}
    
    \noindent\textbf{Resposta:} letra \textbf{\resposta}.
}{%
    \ifdefstring{\modo}{objetivo}{%
        \alternativas
    }{%
        % modo = discursiva → não mostra alternativas
    }
}
\end{exercicioBanco}


    \noindent\rule{\linewidth}{1pt}  % linha horizontal fina
    
    \phantomsection
    \addcontentsline{toc}{section}{Exercício 8}
    \input{../FA_Exercicio8_1}
    \noindent\rule{\linewidth}{1pt}  % linha horizontal fina
    
    \phantomsection
    \addcontentsline{toc}{section}{Exercício 9}
    %%%%%%%%%%%%%%%%%%%%%%%%%%%%%%%%%
%
%       EXERCÍCIO
%
%%%%%%%%%%%%%%%%%%%%%%%%%%%%%%%%%

\ifdefstring{\atividade}{prova}{%
    \ifdefstring{\modo}{objetivo}{%
        \renewcommand{\valorquestao}{\ValorQObj\ ponto}
    }{%
        \renewcommand{\valorquestao}{\ValorQDisc\ pontos}
    }%
}%

\begin{exercicioBanco}[\valorquestao]
Considere a função \(f: \bbR \to \bbR\), \(x \mapsto (3-2a)x + 2\), onde \(a \in \bbR\). Determine a condição sobre o valor \(a\) para que \(f\) seja crescente.

% Define as alternativas
\newcommand{\alternativas}{%
    \begin{center}
        \begin{tabularx}{\textwidth}{XXXXX}
            (a) \(a > 0\). &
            (b) \(a < \dfrac{3}{2}\). &
            (c) \(a = \dfrac{3}{2}\). &
            (d) \(a > \dfrac{3}{2}\). &
            (e) \(a < 3\).
        \end{tabularx}
    \end{center}
}

% Define a resposta correta
\newcommand{\resposta}{B}

% Lógica condicional para exibição
\ifdefstring{\atividade}{lista}{%
    \alternativas
    \vspace{0.5em}
    
    \noindent\textbf{Resposta:} letra \textbf{\resposta}.
}{%
    \ifdefstring{\modo}{objetivo}{%
        \alternativas
    }{%
        % modo = discursiva → não mostra alternativas
    }
}
\end{exercicioBanco}


    \noindent\rule{\linewidth}{1pt}  % linha horizontal fina
    
    \phantomsection
    \addcontentsline{toc}{section}{Exercício 10}
    %%%%%%%%%%%%%%%%%%%%%%%%%%%%%%%%%
%
%       EXERCÍCIO
%
%%%%%%%%%%%%%%%%%%%%%%%%%%%%%%%%%

\ifdefstring{\atividade}{prova}{%
    \ifdefstring{\modo}{objetivo}{%
        \renewcommand{\valorquestao}{\ValorQObj\ ponto}
    }{%
        \renewcommand{\valorquestao}{\ValorQDisc\ pontos}
    }%
}%

\begin{exercicioBanco}[\valorquestao]
Assinale verdadeira (V) ou falso (F) a cada uma das seguintes sentenças.
%
\begin{enumerate}[\(\bullet\)]
    \item A função \(f(x) = 3,5 - 0,4x\) é decrescente.
    %
    \item O zero de \(f(x) = \dfrac{1}{2}x - 2\) é igual a \(2\).
    %
    \item A função \(f(x) = x - 1\) é negativa no intervalo \(]-\infty,1]\).
    %
    \item O coeficiente angular de \(f(x) = 1 - x\) é igual a \(1\).
\end{enumerate}

% Define as alternativas
\newcommand{\alternativas}{%
    \begin{center}
        \begin{tabularx}{\textwidth}{XXXXX}
            (a) (V, F, V, V). &
            (b) (V, V, F, F). &
            (c) (V, F, V, F). &
            (d) (F, V, F, F). &
            (e) (F, V, F, V).
        \end{tabularx}
    \end{center}
}

% Define a resposta correta
\newcommand{\resposta}{C}

% Lógica condicional para exibição
\ifdefstring{\atividade}{lista}{%
    \alternativas
    \vspace{0.5em}
    
    \noindent\textbf{Resposta:} letra \textbf{\resposta}.
}{%
    \ifdefstring{\modo}{objetivo}{%
        \alternativas
    }{%
        % modo = discursiva → não mostra alternativas
    }
}
\end{exercicioBanco}


    \noindent\rule{\linewidth}{1pt}  % linha horizontal fina
    
    \phantomsection
    \addcontentsline{toc}{section}{Exercício 11}
    %%%%%%%%%%%%%%%%%%%%%%%%%%%%%%%%%
%
%       EXERCÍCIO
%
%%%%%%%%%%%%%%%%%%%%%%%%%%%%%%%%%

\ifdefstring{\atividade}{prova}{%
    \ifdefstring{\modo}{objetivo}{%
        \renewcommand{\valorquestao}{\ValorQObj\ ponto}
    }{%
        \renewcommand{\valorquestao}{\ValorQDisc\ pontos}
    }%
}%

\begin{exercicioBanco}[\valorquestao]
Determine o coeficiente angular da função afim \(f: \bbR \to \bbR\) que satisfaz a propriedade de que \(f(x - f(y)) = 1 - x - y\), para qualquer \(x, y \in \bbR\).

% Define as alternativas
\newcommand{\alternativas}{%
    \begin{center}
        \begin{tabularx}{\textwidth}{XXXXX}
            (a) \(0\). &
            (b) \(-2\). &
            (c) \(2\). &
            (d) \(-1\). &
            (e) \(1\).
        \end{tabularx}
    \end{center}
}

% Define a resposta correta
\newcommand{\resposta}{D}

% Lógica condicional para exibição
\ifdefstring{\atividade}{lista}{%
    \alternativas
    \vspace{0.5em}
    
    \noindent\textbf{Resposta:} letra \textbf{\resposta}.
}{%
    \ifdefstring{\modo}{objetivo}{%
        \alternativas
    }{%
        % modo = discursiva → não mostra alternativas
    }
}
\end{exercicioBanco}


    \noindent\rule{\linewidth}{1pt}  % linha horizontal fina
    
    \phantomsection
    \addcontentsline{toc}{section}{Exercício 12}
    %%%%%%%%%%%%%%%%%%%%%%%%%%%%%%%%%
%
%       EXERCÍCIO
%
%%%%%%%%%%%%%%%%%%%%%%%%%%%%%%%%%

\ifdefstring{\atividade}{prova}{%
    \ifdefstring{\modo}{objetivo}{%
        \renewcommand{\valorquestao}{\ValorQObj\ ponto}
    }{%
        \renewcommand{\valorquestao}{\ValorQDisc\ pontos}
    }%
}%

\begin{exercicioBanco}[\valorquestao]
A função que determine o valor a ser pago por uma corrida de táxi é \(f(x) = 3,40 + 2,50x\), sendo \(x\) a distância percorrida em km. Determine o valor a ser pago por uma corrida de 10km.

% Define as alternativas
\newcommand{\alternativas}{%
    \begin{center}
        \begin{tabularx}{\textwidth}{XXXXX}
            (a) R\$ \(25{,}00\). &
            (b) R\$ \(28{,}40\). &
            (c) R\$ \(34{,}00\). &
            (d) R\$ \(30{,}90\). &
            (e) R\$ \(27{,}50\).
        \end{tabularx}
    \end{center}
}

% Resposta correta
\newcommand{\resposta}{B}

% Lógica condicional para exibição
\ifdefstring{\atividade}{lista}{%
    \alternativas
    \vspace{0.5em}
    
    \noindent\textbf{Resposta:} letra \textbf{\resposta}.
}{%
    \ifdefstring{\modo}{objetivo}{%
        \alternativas
    }{%
        % modo = discursiva → não mostra alternativas
    }
}
\end{exercicioBanco}


    \noindent\rule{\linewidth}{1pt}  % linha horizontal fina
    
    \phantomsection
    \addcontentsline{toc}{section}{Exercício 13}
    %%%%%%%%%%%%%%%%%%%%%%%%%%%%%%%%%
%
%       EXERCÍCIO
%
%%%%%%%%%%%%%%%%%%%%%%%%%%%%%%%%%

\ifdefstring{\atividade}{prova}{%
    \ifdefstring{\modo}{objetivo}{%
        \renewcommand{\valorquestao}{\ValorQObj\ ponto}
    }{%
        \renewcommand{\valorquestao}{\ValorQDisc\ pontos}
    }%
}%

\begin{exercicioBanco}[\valorquestao]
Uma cisterna de 6000 litros foi esvaziada em um período de 3h. Na primeira hora foi utilizada apenas uma bomba, mas nas duas horas seguintes, a fim de reduzir o tempo de esvaziamento, outra bomba foi ligada junto com a primeira. O gráfico, formado por dois segmentos de reta, mostra o volume de água presente na cisterna, em função do tempo

%
\begin{center}
    \begin{tikzpicture}[scale=0.9]
        % Eixos
        \draw[->] (0,0) -- (3.5,0) node[below right] {\small{Tempo (horas)}};
        \draw[->] (0,0) -- (0,3.5) node[above left] {\small{Volume (litros)}};
              
        % Pontos escolhidos
        \draw[dashed] (1,0) -- (1,2.5) -- (0,2.5);
             
        % Marcando os pontos
        \filldraw (0,3) circle (1pt) node[left] {\small\(6000\)};
        \filldraw (1,2.5) circle (1pt) node[xshift=-41pt] {\small\(5000\)};
        \filldraw (1,0) circle (1pt) node[below] {\small\(1\)};
        \filldraw (3,0) circle (1pt) node[below] {\small\(3\)};
        
        \draw[thick, blue] (0,3) -- (1,2.5) -- (3,0);
    \end{tikzpicture}
\end{center}
%

%
Encontre a vazão, em litro por hora, da bomba que foi ligada no início da segunda hora.
%

% Define as alternativas
\newcommand{\alternativas}{%
    \begin{center}
        \begin{tabularx}{\textwidth}{XXXXX}
            (a) \(500\). &
            (b) \(1500\). &
            (c) \(2000\). &
            (d) \(2500\). &
            (e) \(3000\).
        \end{tabularx}
    \end{center}
}

% Resposta correta
\newcommand{\resposta}{B}

% Lógica condicional para exibição
\ifdefstring{\atividade}{lista}{%
    \alternativas
    \vspace{0.5em}
    
    \noindent\textbf{Resposta:} letra \textbf{\resposta}.
}{%
    \ifdefstring{\modo}{objetivo}{%
        \alternativas
    }{%
        % modo = discursiva → não mostra alternativas
    }
}
\end{exercicioBanco}


    \noindent\rule{\linewidth}{1pt}  % linha horizontal fina
    
    \phantomsection
    \addcontentsline{toc}{section}{Exercício 14}
    %%%%%%%%%%%%%%%%%%%%%%%%%%%%%%%%%
%
%       EXERCÍCIO
%
%%%%%%%%%%%%%%%%%%%%%%%%%%%%%%%%%

\ifdefstring{\atividade}{prova}{%
    \ifdefstring{\modo}{objetivo}{%
        \renewcommand{\valorquestao}{\ValorQObj\ ponto}
    }{%
        \renewcommand{\valorquestao}{\ValorQDisc\ pontos}
    }%
}%

\begin{exercicioBanco}[\valorquestao]
Dois objetos metálicos, ambos com temperatura inicial igual a \(10\,^\circ\text{C}\), são aquecidos. Suas temperaturas, \(y_{1}\) e \(y_{2}\), em função do tempo \(t\), em segundo, estão representadas no plano cartesiano pelas semirretas com origem no ponto \(A = (0,10)\) e que passam, respectivamente, pelos pontos \(B = (20, 50)\) e \(C = (40, 30)\). Sabe-se que, em determinado intervalo de tempo, a temperatura \(y_{1}\) aumentou \(20\,^\circ\text{C}\).  

%
\begin{center}
    \begin{tikzpicture}[scale=0.5]
        % Eixos
        \draw[->] (0,0) -- (5.5,0) node[below right] {\small{t (s)}};
        \draw[->] (0,0) -- (0,5.5) node[above left] {\small{y (\(^\circ\text{C}\))}};
        
        % Gráfico
        \draw[domain=0:2.1, smooth, variable=\x, blue, thick] plot ({\x}, {2*\x+1}); \node[blue, above left] at (1.7,3.5) {\small\(y_{1}\)};
        \draw[domain=0:4.5, smooth, variable=\x, red, thick] plot ({\x}, {(1/2)*\x+1}); \node[red, below] at (3.5,2.3) {\small\(y_{2}\)};
             
        % Marcando os pontos
        \filldraw (4,3) circle (2pt) node[above left] {\small\(C\)};
        \filldraw (2,5) circle (2pt) node[below right] {\small\(B\)};
        \filldraw (0,1) circle (2pt) node[left] {\small\(A\)};
    \end{tikzpicture}
\end{center}
%

Nesse mesmo intervalo de tempo, determine quanto aumentou a temperatura \(y_{2}\), em grau Celsius.

% Define as alternativas
\newcommand{\alternativas}{%
    \begin{center}
        \begin{tabularx}{\textwidth}{XXXXX}
            (a) \(5\). &
            (b) \(10\). &
            (c) \(15\). &
            (d) \(20\). &
            (e) \(25\).
        \end{tabularx}
    \end{center}
}

% Resposta correta
\newcommand{\resposta}{A}

% Lógica condicional para exibição
\ifdefstring{\atividade}{lista}{%
    \alternativas
    \vspace{0.5em}
    
    \noindent\textbf{Resposta:} letra \textbf{\resposta}.
}{%
    \ifdefstring{\modo}{objetivo}{%
        \alternativas
    }{%
        % modo = discursiva → não mostra alternativas
    }
}
\end{exercicioBanco}


    \noindent\rule{\linewidth}{1pt}  % linha horizontal fina
    
    \phantomsection
    \addcontentsline{toc}{section}{Exercício 15}
    %%%%%%%%%%%%%%%%%%%%%%%%%%%%%%%%%
%
%       EXERCÍCIO
%
%%%%%%%%%%%%%%%%%%%%%%%%%%%%%%%%%

\ifdefstring{\atividade}{prova}{%
    \ifdefstring{\modo}{objetivo}{%
        \renewcommand{\valorquestao}{\ValorQObj\ ponto}
    }{%
        \renewcommand{\valorquestao}{\ValorQDisc\ pontos}
    }%
}%

\begin{exercicioBanco}[\valorquestao]
Uma empresa produz mochilas escolares sob encomenda. Essa empresa tem um custo total de produção, composto por um custo fixo, que não depende do número de mochilas, mais um custo variável, que é proporcional ao número de mochilas produzidas. O custo total cresce de forma linear, e a tabela apresenta esse custo para três quantidades de mochilas produzidas. 
%
\begin{center}
    \begin{tabular}{|C{200pt}|C{50pt}|C{50pt}|C{50pt}|}
        \hline
        Quantidade de mochilas & 30 & 50 & 100 \\
        \hline
        Custo total (R\$) & 1050,00 & 1650,00 & 3150,00 \\
        \hline \hline
    \end{tabular}
\end{center}
%
Responda qual será o custo total, em real, para a produção de 80 mochilas.

% Define as alternativas
\newcommand{\alternativas}{%
    \begin{center}
        \begin{tabularx}{\textwidth}{XXXXX}
            (a) R\$ \(2400{,}00\). &
            (b) R\$ \(2550{,}00\). &
            (c) R\$ \(2700{,}00\). &
            (d) R\$ \(2850{,}00\). &
            (e) R\$ \(3000{,}00\).
        \end{tabularx}
    \end{center}
}

% Resposta correta
\newcommand{\resposta}{C}

% Lógica condicional para exibição
\ifdefstring{\atividade}{lista}{%
    \alternativas
    \vspace{0.5em}
    
    \noindent\textbf{Resposta:} letra \textbf{\resposta}.
}{%
    \ifdefstring{\modo}{objetivo}{%
        \alternativas
    }{%
        % modo = discursiva → não mostra alternativas
    }
}
\end{exercicioBanco}


    \noindent\rule{\linewidth}{1pt}  % linha horizontal fina
    
    \phantomsection
    \addcontentsline{toc}{section}{Exercício 16}
    %%%%%%%%%%%%%%%%%%%%%%%%%%%%%%%%%
%
%       EXERCÍCIO
%
%%%%%%%%%%%%%%%%%%%%%%%%%%%%%%%%%

\ifdefstring{\atividade}{prova}{%
    \ifdefstring{\modo}{objetivo}{%
        \renewcommand{\valorquestao}{\ValorQObj\ ponto}
    }{%
        \renewcommand{\valorquestao}{\ValorQDisc\ pontos}
    }%
}%

\begin{exercicioBanco}[\valorquestao]
Os gráficos abaixo representam as funções receita mensal \(R(x)\) e custo mensal \(C(x)\) de um produto fabricado por uma empresa, em que \(x\) é a quantidade produzida e vendida. Responda qual é o lucro obtido ao se produzir
e vender \(1350\) unidades por mês.

\begin{center}
    \begin{tikzpicture}[scale=0.5]

        % Eixos
        \draw[->] (0,0) -- (8,0) node[below right] {\small Quantidade};
        \draw[->] (0,0) -- (0,6) node[above left] {\small Receita e Custo};

        % R(x) = 15x (normalizado para caber melhor)
        \draw[domain=0:7, smooth, variable=\x, blue, thick] 
            plot ({\x}, {0.8*\x});
        \node[blue, above] at (6,5.2) {\small \(R(x)\)};

        % C(x) = 10x + 5000 (também reescalado)
        \draw[domain=0:7, smooth, variable=\x, red, thick] 
            plot ({\x}, {0.5*\x + 1.5});
        \node[red, right] at (6,4) {\small \(C(x)\)};

        % Ponto de interseção (aproximado)
        \filldraw (5,4) circle (1.2pt);
        
        \draw[dashed] (5,0) -- (5,4) -- (0,4);
        \node[below] at (5,0) {\small \(1000\)};
        \node[left] at (0,4) {\small \(15000\)};
        \node[below left] at (0,0) {\small \(0\)};
        \node[left] at (0,1.5) {\small \(5000\)};

        % Marcas nos eixos (opcional)
        \foreach \x in {1,2,3,4,5,6,7}
            \draw (\x,0) -- (\x,-0.1);

        \foreach \y in {1,2,3,4,5}
            \draw (0,\y) -- (-0.1,\y);

    \end{tikzpicture}
\end{center}

% Define as alternativas
\newcommand{\alternativas}{%
    \begin{center}
        \begin{tabularx}{\textwidth}{XXXXX}
            (a) R\$ \(6750{,}00\). &
            (b) R\$ \(2500{,}00\). &
            (c) R\$ \(2000{,}00\). &
            (d) R\$ \(1750{,}00\). &
            (e) R\$ \(1500{,}00\).
        \end{tabularx}
    \end{center}
}

% Resposta correta
\newcommand{\resposta}{D}

% Lógica condicional para exibição
\ifdefstring{\atividade}{lista}{%
    \alternativas
    \vspace{0.5em}
    
    \noindent\textbf{Resposta:} letra \textbf{\resposta}.
}{%
    \ifdefstring{\modo}{objetivo}{%
        \alternativas
    }{%
        % modo = discursiva → não mostra alternativas
    }
}
\end{exercicioBanco}


    \noindent\rule{\linewidth}{1pt}  % linha horizontal fina
    
    \phantomsection
    \addcontentsline{toc}{section}{Exercício 17}
    %%%%%%%%%%%%%%%%%%%%%%%%%%%%%%%%%
%
%       EXERCÍCIO
%
%%%%%%%%%%%%%%%%%%%%%%%%%%%%%%%%%

\ifdefstring{\atividade}{prova}{%
    \ifdefstring{\modo}{objetivo}{%
        \renewcommand{\valorquestao}{\ValorQObj\ ponto}
    }{%
        \renewcommand{\valorquestao}{\ValorQDisc\ pontos}
    }%
}%

\begin{exercicioBanco}[\valorquestao]
Considerando um intervalo de tempo de 10 anos a partir de hoje, o valor de uma máquina deprecia linearmente com o tempo, isto é, o valor da máquina \(y\), em função do tempo \(x\), é dado por uma função afim \(y = ax + b\). Se o valor da máquina daqui a dois anos for R\$ \(6400{,}00\), e seu valor daqui a cinco anos e meio for R\$ \(4300{,}00\), determine qual será o seu valor daqui a sete anos.

% Define as alternativas
\newcommand{\alternativas}{%
    \begin{center}
        \begin{tabularx}{\textwidth}{XXXXX}
            (a) R\$ \(2800{,}00\). &
            (b) R\$ \(3000{,}00\). &
            (c) R\$ \(3200{,}00\). &
            (d) R\$ \(3400{,}00\). &
            (e) R\$ \(3600{,}00\).
        \end{tabularx}
    \end{center}
}

% Resposta correta
\newcommand{\resposta}{D}

% Lógica condicional para exibição
\ifdefstring{\atividade}{lista}{%
    \alternativas
    \vspace{0.5em}
    
    \noindent\textbf{Resposta:} letra \textbf{\resposta}.
}{%
    \ifdefstring{\modo}{objetivo}{%
        \alternativas
    }{%
        % modo = discursiva → não mostra alternativas
    }
}
\end{exercicioBanco}


    \noindent\rule{\linewidth}{1pt}  % linha horizontal fina
    
    \phantomsection
    \addcontentsline{toc}{section}{Exercício 18}
    %%%%%%%%%%%%%%%%%%%%%%%%%%%%%%%%%
%
%       EXERCÍCIO
%
%%%%%%%%%%%%%%%%%%%%%%%%%%%%%%%%%

\ifdefstring{\atividade}{prova}{%
    \ifdefstring{\modo}{objetivo}{%
        \renewcommand{\valorquestao}{\ValorQObj\ ponto}
    }{%
        \renewcommand{\valorquestao}{\ValorQDisc\ pontos}
    }%
}%

\begin{exercicioBanco}[\valorquestao]
Para garantir a segurança de um grande evento público que terá início às 4h da tarde, um organizador precisa monitorar a quantidade de pessoas presentes em cada instante. Para cada 2.000 pessoas, faz-se necessária a presença de um policial. Além disso, estima-se uma densidade de quatro pessoas por metro quadrado de área de terreno ocupado.

Às 10h da manhã, o organizador verifica que a área de terreno já ocupada equivale a um quadrado com lados medindo \(500\) m. Porém, nas horas seguintes, espera-se que o público aumente a uma taxa de \(120.000\) pessoas por hora até o início do evento, quando não será mais permitida a entrada de público.

Responda quantos policiais serão necessários no início do evento para garantir a segurança?

% Define as alternativas
\newcommand{\alternativas}{%
    \begin{center}
        \begin{tabularx}{\textwidth}{XXXXX}
            (a) \(720\). &
            (b) \(780\). &
            (c) \(800\). &
            (d) \(820\). &
            (e) \(860\).
        \end{tabularx}
    \end{center}
}

% Resposta correta
\newcommand{\resposta}{E}

% Lógica condicional para exibição
\ifdefstring{\atividade}{lista}{%
    \alternativas
    \vspace{0.5em}
    
    \noindent\textbf{Resposta:} letra \textbf{\resposta}.
}{%
    \ifdefstring{\modo}{objetivo}{%
        \alternativas
    }{%
        % modo = discursiva → não mostra alternativas
    }
}
\end{exercicioBanco}


    \noindent\rule{\linewidth}{1pt}  % linha horizontal fina
    
\end{document}
%*************** FINAL DO DOCUMENTO ***************
